\section{Inline Visual Route: Evidence and Replay}

These figures teach how numeric and computational evidence enters the public claim surface.


\subsection{Replay row anatomy}

This visual anchor names the objects, review direction, and formula or numeric boundary that keep the figure tied to the argument.


\begin{figure}[htbp]

\centering

\begin{tikzpicture}[>=Latex,node distance=2.4cm]

\node[draw,rounded corners,fill=blue!7,text width=0.28\textwidth,align=center,minimum height=1.0cm] (a) {Formula};

\node[draw,rounded corners,fill=green!7,text width=0.28\textwidth,align=center,minimum height=1.0cm,right=of a] (b) {Observed value};

\draw[->,line width=0.8pt] (a) -- node[above,align=center] {residual} (b);

\node[draw,rounded corners,fill=orange!8,text width=0.68\textwidth,align=center,below=0.9cm of $(a)!0.5!(b)$] (c) {review question: support class, boundary, falsifier};

\draw[->,dashed] (c.north) -- ($(a)!0.5!(b)$);

\end{tikzpicture}

\caption{Replay row anatomy. Shows formula, observation, and residual as one evidence row. The figure is placed inline in the manuscript; complete machine evidence remains in the evidence package.}

\label{fig:r005-inline-28c_oc133_inline_figures_evidence_route-1}

\end{figure}

\subsection{Target-blind split}

This visual anchor names the objects, review direction, and formula or numeric boundary that keep the figure tied to the argument.


\begin{figure}[htbp]

\centering

\begin{tikzpicture}[>=Latex,node distance=2.4cm]

\node[draw,rounded corners,fill=blue!7,text width=0.28\textwidth,align=center,minimum height=1.0cm] (a) {Held target};

\node[draw,rounded corners,fill=green!7,text width=0.28\textwidth,align=center,minimum height=1.0cm,right=of a] (b) {Unlocked result};

\draw[->,line width=0.8pt] (a) -- node[above,align=center] {split policy} (b);

\node[draw,rounded corners,fill=orange!8,text width=0.68\textwidth,align=center,below=0.9cm of $(a)!0.5!(b)$] (c) {review question: support class, boundary, falsifier};

\draw[->,dashed] (c.north) -- ($(a)!0.5!(b)$);

\end{tikzpicture}

\caption{Target-blind split. Shows why target-blind replay is stronger than archive lookup. The figure is placed inline in the manuscript; complete machine evidence remains in the evidence package.}

\label{fig:r005-inline-28c_oc133_inline_figures_evidence_route-2}

\end{figure}

\subsection{Comparator lane}

This visual anchor names the objects, review direction, and formula or numeric boundary that keep the figure tied to the argument.


\begin{figure}[htbp]

\centering

\begin{tikzpicture}[>=Latex,node distance=2.4cm]

\node[draw,rounded corners,fill=blue!7,text width=0.28\textwidth,align=center,minimum height=1.0cm] (a) {OC row};

\node[draw,rounded corners,fill=green!7,text width=0.28\textwidth,align=center,minimum height=1.0cm,right=of a] (b) {Baseline};

\draw[->,line width=0.8pt] (a) -- node[above,align=center] {delta} (b);

\node[draw,rounded corners,fill=orange!8,text width=0.68\textwidth,align=center,below=0.9cm of $(a)!0.5!(b)$] (c) {review question: support class, boundary, falsifier};

\draw[->,dashed] (c.north) -- ($(a)!0.5!(b)$);

\end{tikzpicture}

\caption{Comparator lane. Shows comparison as residual-delta accounting rather than priority rhetoric. The figure is placed inline in the manuscript; complete machine evidence remains in the evidence package.}

\label{fig:r005-inline-28c_oc133_inline_figures_evidence_route-3}

\end{figure}

\subsection{Uncertainty route}

This visual anchor names the objects, review direction, and formula or numeric boundary that keep the figure tied to the argument.


\begin{figure}[htbp]

\centering

\begin{tikzpicture}[>=Latex,node distance=2.4cm]

\node[draw,rounded corners,fill=blue!7,text width=0.28\textwidth,align=center,minimum height=1.0cm] (a) {Prediction};

\node[draw,rounded corners,fill=green!7,text width=0.28\textwidth,align=center,minimum height=1.0cm,right=of a] (b) {Interval};

\draw[->,line width=0.8pt] (a) -- node[above,align=center] {error bar} (b);

\node[draw,rounded corners,fill=orange!8,text width=0.68\textwidth,align=center,below=0.9cm of $(a)!0.5!(b)$] (c) {review question: support class, boundary, falsifier};

\draw[->,dashed] (c.north) -- ($(a)!0.5!(b)$);

\end{tikzpicture}

\caption{Uncertainty route. Shows that numbers require uncertainty or residual interpretation. The figure is placed inline in the manuscript; complete machine evidence remains in the evidence package.}

\label{fig:r005-inline-28c_oc133_inline_figures_evidence_route-4}

\end{figure}

\subsection{Falsifier row}

This visual anchor names the objects, review direction, and formula or numeric boundary that keep the figure tied to the argument.


\begin{figure}[htbp]

\centering

\begin{tikzpicture}[>=Latex,node distance=2.4cm]

\node[draw,rounded corners,fill=blue!7,text width=0.28\textwidth,align=center,minimum height=1.0cm] (a) {Promoted claim};

\node[draw,rounded corners,fill=green!7,text width=0.28\textwidth,align=center,minimum height=1.0cm,right=of a] (b) {Failing observation};

\draw[->,line width=0.8pt] (a) -- node[above,align=center] {reopen} (b);

\node[draw,rounded corners,fill=orange!8,text width=0.68\textwidth,align=center,below=0.9cm of $(a)!0.5!(b)$] (c) {review question: support class, boundary, falsifier};

\draw[->,dashed] (c.north) -- ($(a)!0.5!(b)$);

\end{tikzpicture}

\caption{Falsifier row. Shows how an empirical claim can be reopened. The figure is placed inline in the manuscript; complete machine evidence remains in the evidence package.}

\label{fig:r005-inline-28c_oc133_inline_figures_evidence_route-5}

\end{figure}

\subsection{Evidence package link}

This visual anchor names the objects, review direction, and formula or numeric boundary that keep the figure tied to the argument.


\begin{figure}[htbp]

\centering

\begin{tikzpicture}[>=Latex,node distance=2.4cm]

\node[draw,rounded corners,fill=blue!7,text width=0.28\textwidth,align=center,minimum height=1.0cm] (a) {Prose claim};

\node[draw,rounded corners,fill=green!7,text width=0.28\textwidth,align=center,minimum height=1.0cm,right=of a] (b) {Machine row};

\draw[->,line width=0.8pt] (a) -- node[above,align=center] {hash} (b);

\node[draw,rounded corners,fill=orange!8,text width=0.68\textwidth,align=center,below=0.9cm of $(a)!0.5!(b)$] (c) {review question: support class, boundary, falsifier};

\draw[->,dashed] (c.north) -- ($(a)!0.5!(b)$);

\end{tikzpicture}

\caption{Evidence package link. Shows why exact rows live in the package while interpretation stays in prose. The figure is placed inline in the manuscript; complete machine evidence remains in the evidence package.}

\label{fig:r005-inline-28c_oc133_inline_figures_evidence_route-6}

\end{figure}
